% GENERATED FILE. Edit canonical lessons, then run npm run generate:books.
% canonical-source-hash: fnv1a64:c47f93a7953f17ae
% canonical-lessons: JA-C142-chuui, JA-C142-kinen, JA-C142-uketsuke, JA-C142-eigyouchuu, JA-C142-junbichuu, JA-R142-first-pass-my-day-no-and-the-past, JA-R142-second-pass-the-future-and-notices

\chapter{Notices: Caution, No Smoking, Reception, Open}
\label{ch:ja-notices-caution-no-smoking-reception-open}
\hlchaptermodality{142}

\begin{chapteropening}
\textbf{By the end of this chapter:} \emph{I can read short notices and door signs, such as caution, no smoking, reception, open and not yet open, and act on what they ask of me.}

Read caution, no smoking, reception, open and not yet open, then read a three-line notice on a clinic door and say what you do; say next week, next month, next year, a plan and probably.
\end{chapteropening}

\section[\texorpdfstring{\ja{ちゅうい} (chūi)}{ちゅうい (chūi)}]{\ja{ちゅうい} (chūi) — caution, take care (on a sign)}

\label{lesson:JA-C142-chuui}



\begin{quote}
Four recalls before the new word.

\begin{itemize}
\raggedright
  \item \emph{Recall:} say \emph{probably} — \textbf{R1}, one lesson back
  \item \emph{Recall:} say \emph{next week} — \textbf{R2}, five lessons back
  \item \emph{Recall:} say \emph{I work}, politely — \textbf{R3}, twenty lessons back
  \item \emph{Write it:} \textbf{\ja{る}} from memory — \textbf{R4}, eighty lessons back
\end{itemize}
\end{quote}

\subsection*{You'll want to know: \ja{ちゅうい}}
\textbf{\ja{ちゅうい}} — \emph{chūi} — \textquotedblleft{}caution\textquotedblright{}, \textquotedblleft{}take care\textquotedblright{}. It heads warning signs and notices. The small \textbf{\ja{ゅ}} joins \textbf{\ja{ち}} into one beat, \emph{chu}, and \textbf{\ja{う}} holds it long: \emph{chū}, then \emph{i}. On a real sign this word is usually printed in kanji, which this book has not reached yet; the reading is the same.

\begin{quote}
\textbf{\ja{くるまに} \ja{ちゅうい}} — \emph{kuruma ni chūi} — \textquotedblleft{}Watch out for cars.\textquotedblright{}
\end{quote}

Read a notice like this and act on it: stop, look, and cross with care.

\subsection*{Your turn}
\begin{itemize}
\raggedright
  \item \emph{Say it:} \emph{chūi}
  \item \emph{Read it:} \textbf{\ja{くるまに} \ja{ちゅうい}} — and say what it asks you to do
  \item \emph{Recall:} say \emph{a car}, then read \emph{watch out for cars} aloud
\end{itemize}

\begin{tcolorbox}[breakable,colback=teal!4,colframe=teal!35!black,title={Before you move on}]
What does \ja{ちゅうい} on a sign mean? (\textquotedblleft{}Caution\textquotedblright{}, take care.) Say it once more.
\end{tcolorbox}
\section[\texorpdfstring{\ja{きんえん} (kin'en)}{きんえん (kin'en)}]{\ja{きんえん} (kin'en) — no smoking (on a sign)}

\label{lesson:JA-C142-kinen}



\begin{quote}
Three recalls before the new word.

\begin{itemize}
\raggedright
  \item \emph{Recall:} read \emph{caution} on a sign — \textbf{R1}, one lesson back
  \item \emph{Recall:} say \emph{next month} — \textbf{R2}, five lessons back
  \item \emph{Recall:} say \emph{I walk}, politely — \textbf{R3}, twenty lessons back
\end{itemize}
\end{quote}

\subsection*{You'll want to know: \ja{きんえん}}
\textbf{\ja{きんえん}} — \emph{kin'en} — \textquotedblleft{}no smoking\textquotedblright{}. The apostrophe in the romanization keeps \emph{kin} and \emph{en} apart: four beats, \emph{ki–n–e–n}. On a real sign this word is usually printed in kanji, which this book has not reached yet; the reading is the same.

\begin{quote}
\textbf{\ja{きんえん}}
\end{quote}

\begin{quote}
\textbf{\ja{ここで} \ja{たばこを} \ja{すわないで} \ja{ください。}} — \emph{koko de tabako o suwanaide kudasai} — \textquotedblleft{}Please do not smoke here.\textquotedblright{}
\end{quote}

\textbf{\ja{すう}}, \emph{suu}, the verb you know for breathing in, also means to smoke. Its negative stem \textbf{\ja{すわ}} with \textbf{\ja{ないで} \ja{ください}} asks you politely not to do it.

\subsection*{Your turn}
\begin{itemize}
\raggedright
  \item \emph{Say it:} \emph{kin'en}
  \item \emph{Read it:} \textbf{\ja{ここで} \ja{たばこを} \ja{すわないで} \ja{ください}} — and say what it asks
  \item \emph{Recall:} read \emph{caution}, then \emph{no smoking}, as they appear on signs
\end{itemize}

\begin{tcolorbox}[breakable,colback=teal!4,colframe=teal!35!black,title={Before you move on}]
What does \ja{きんえん} on a sign mean? (\textquotedblleft{}No smoking\textquotedblright{}.) Say it once more.
\end{tcolorbox}
\section[\texorpdfstring{\ja{うけつけ} (uketsuke)}{うけつけ (uketsuke)}]{\ja{うけつけ} (uketsuke) — a reception desk, reception}

\label{lesson:JA-C142-uketsuke}



\begin{quote}
Four recalls before the new word.

\begin{itemize}
\raggedright
  \item \emph{Recall:} read \emph{no smoking} on a sign — \textbf{R1}, one lesson back
  \item \emph{Recall:} say \emph{next year} — \textbf{R2}, five lessons back
  \item \emph{Recall:} say \emph{I go home}, politely — \textbf{R3}, twenty lessons back
  \item \emph{Recall:} say \emph{a station} — \textbf{R4}, eighty lessons back
\end{itemize}
\end{quote}

\subsection*{You'll want to know: \ja{うけつけ}}
\textbf{\ja{うけつけ}} — \emph{uketsuke} — \textquotedblleft{}reception\textquotedblright{}, \textquotedblleft{}a reception desk\textquotedblright{}: where you check in at a hospital, an office or a school. On a real sign this word is usually printed in kanji, which this book has not reached yet; the reading is the same.

\begin{quote}
\textbf{\ja{うけつけで} \ja{なまえを} \ja{かいて} \ja{ください。}} — \emph{uketsuke de namae o kaite kudasai} — \textquotedblleft{}Please write your name at reception.\textquotedblright{}
\end{quote}

\textbf{\ja{かく}}, \emph{kaku}, \textquotedblleft{}to write\textquotedblright{}, becomes \textbf{\ja{かいて}} before \textbf{\ja{ください}}. The notice tells you what to do: go to the desk, and write your name.

\subsection*{Your turn}
\begin{itemize}
\raggedright
  \item \emph{Say it:} \emph{uketsuke}
  \item \emph{Read it:} \textbf{\ja{うけつけで} \ja{なまえを} \ja{かいて} \ja{ください}} — and say what you do first
  \item \emph{Recall:} say \emph{a hospital}, then where you go first inside it
\end{itemize}

\begin{tcolorbox}[breakable,colback=teal!4,colframe=teal!35!black,title={Before you move on}]
What is \ja{うけつけ}? (\textquotedblleft{}Reception\textquotedblright{}, the desk where you check in.) Say it once more.
\end{tcolorbox}
\section[\texorpdfstring{\ja{えいぎょうちゅう} (eigyōchū)}{えいぎょうちゅう (eigyōchū)}]{\ja{えいぎょうちゅう} (eigyōchū) — open, open for business (on a shop door)}

\label{lesson:JA-C142-eigyouchuu}



\begin{quote}
Four recalls before the new word.

\begin{itemize}
\raggedright
  \item \emph{Recall:} read \emph{reception} on a sign — \textbf{R1}, one lesson back
  \item \emph{Recall:} say \emph{a plan}, what you mean to do — \textbf{R2}, five lessons back
  \item \emph{Recall:} say \emph{I go to bed}, politely — \textbf{R3}, twenty lessons back
  \item \emph{Write it:} \textbf{\ja{き}} from memory — \textbf{R4}, eighty lessons back
\end{itemize}
\end{quote}

\subsection*{You'll want to know: \ja{えいぎょうちゅう}}
\textbf{\ja{えいぎょうちゅう}} — \emph{eigyōchū} — \textquotedblleft{}open\textquotedblright{}, \textquotedblleft{}open for business\textquotedblright{}, on a small sign hung on a shop or restaurant door. \textbf{\ja{ちゅう}} at the end means \textquotedblleft{}in the middle of\textquotedblright{}: the shop is in the middle of doing business. Eight signs make six beats, \emph{e–i–gyo–o–chu–u}: each small sign joins the one before it, and each \textbf{\ja{う}} holds the sound before it long. On a real sign this word is usually printed in kanji, which this book has not reached yet; the reading is the same.

\emph{Read it:} it and act on it — the door says \textbf{\ja{えいぎょうちゅう}}, so you can go in

\subsection*{Your turn}
\begin{itemize}
\raggedright
  \item \emph{Say it:} \emph{eigyōchū}
  \item \emph{Read it:} \textbf{\ja{えいぎょうちゅう}} on a door — and say whether you can go in
  \item \emph{Recall:} read \emph{reception}, then \emph{open}
\end{itemize}

\begin{tcolorbox}[breakable,colback=teal!4,colframe=teal!35!black,title={Before you move on}]
What does \ja{えいぎょうちゅう} on a door mean? (\textquotedblleft{}Open\textquotedblright{}, for business.) Say it once more.
\end{tcolorbox}
\section[\texorpdfstring{\ja{じゅんびちゅう} (junbichū)}{じゅんびちゅう (junbichū)}]{\ja{じゅんびちゅう} (junbichū) — getting ready, not yet open (on a shop door)}

\label{lesson:JA-C142-junbichuu}



\begin{quote}
Four recalls before the new word.

\begin{itemize}
\raggedright
  \item \emph{Recall:} read \emph{open} on a shop door — \textbf{R1}, one lesson back
  \item \emph{Recall:} say \emph{probably} — \textbf{R2}, five lessons back
  \item \emph{Recall:} ask \emph{is it ...?}, politely — \textbf{R3}, twenty lessons back
  \item \emph{Recall:} say \emph{yesterday} — \textbf{R4}, eighty lessons back
\end{itemize}
\end{quote}

\subsection*{You'll want to know: \ja{じゅんびちゅう}}
\textbf{\ja{じゅんびちゅう}} — \emph{junbichū} — \textquotedblleft{}getting ready\textquotedblright{}, \textquotedblleft{}not yet open\textquotedblright{}: the other side of the same door sign. \textbf{\ja{じゅんび}} is preparation, and \textbf{\ja{ちゅう}} again means \textquotedblleft{}in the middle of\textquotedblright{}. On a real sign this word is usually printed in kanji, which this book has not reached yet; the reading is the same.

\emph{Read it:} it and act on it — the door says \textbf{\ja{じゅんびちゅう}}, so the shop is not open yet

Wait, or come back later.

\begin{quote}
\textbf{\ja{いまは} \ja{じゅんびちゅうです。}} — \emph{ima wa junbichū desu} — \textquotedblleft{}We are getting ready now; not open yet.\textquotedblright{}
\end{quote}

\subsection*{Your turn}
\begin{itemize}
\raggedright
  \item \emph{Say it:} \emph{junbichū}
  \item \emph{Read it:} \textbf{\ja{じゅんびちゅう}} on a door — and say what you do
  \item \emph{Recall:} read \emph{open}, then \emph{not yet open}, and say which one lets you in
\end{itemize}

\begin{tcolorbox}[breakable,colback=teal!4,colframe=teal!35!black,title={Before you move on}]
What does \ja{じゅんびちゅう} on a door mean? (\textquotedblleft{}Getting ready\textquotedblright{}, not yet open.) Say it once more.
\end{tcolorbox}
\section[(dialogue)]{My day, saying no, and the past — a first pass}

\label{lesson:JA-R142-first-pass-my-day-no-and-the-past}



\begin{quote}
Four recalls, then the review.

\begin{itemize}
\raggedright
  \item \emph{Recall:} read \emph{getting ready, not yet open} on a shop door — \textbf{R1}, one lesson back
  \item \emph{Recall:} read \emph{caution} on a sign — \textbf{R2}, five lessons back
  \item \emph{Recall:} answer \emph{no, that is not right} — \textbf{R3}, twenty lessons back
  \item \emph{Recall:} say \emph{a pond} — \textbf{R4}, eighty lessons back
\end{itemize}
\end{quote}

\subsection*{Your turn}
Say each one in Japanese.

Say these aloud:

\begin{itemize}
\raggedright
  \item I get up, I work, I walk, I go home, I go to bed — politely
  \item is it ...?, no, that is not right, is not (plainly)
  \item nothing, nobody, nowhere — each with a negative verb
  \item last week, last month, last year, a moment ago, it is over
\end{itemize}

Then say three sentences: what you did last week (\textbf{\ja{ました}}), what you did not do last year (\textbf{\ja{ませんでした}}), and a question about tomorrow (\textbf{\ja{ますか}}).

\begin{tcolorbox}[breakable,colback=teal!4,colframe=teal!35!black,title={Before you move on}]
I get up? (\textbf{\ja{おきます}}, \emph{okimasu}.) I work? (\textbf{\ja{はたらきます}}, \emph{hatarakimasu}.) I walk? (\textbf{\ja{あるきます}}, \emph{arukimasu}.) I go home? (\textbf{\ja{かえります}}, \emph{kaerimasu}.) I go to bed? (\textbf{\ja{ねます}}, \emph{nemasu}.) Is it ...? (\textbf{\ja{ですか}}, \emph{desu ka}.) No, that is not right? (\textbf{\ja{ちがいます}}, \emph{chigaimasu}.) Is not? (\textbf{\ja{じゃない}}, \emph{ja nai}.) Nothing? (\textbf{\ja{なにも}}, \emph{nani mo}.) Nobody? (\textbf{\ja{だれも}}, \emph{dare mo}.) Nowhere? (\textbf{\ja{どこにも}}, \emph{doko ni mo}.) Last week? (\textbf{\ja{せんしゅう}}, \emph{senshū}.) Last month? (\textbf{\ja{せんげつ}}, \emph{sengetsu}.) Last year? (\textbf{\ja{きょねん}}, \emph{kyonen}.) A moment ago? (\textbf{\ja{さっき}}, \emph{sakki}.) It is over? (\textbf{\ja{おわりました}}, \emph{owarimashita}.)
\end{tcolorbox}
\section[(dialogue)]{The future and notices — a second pass}

\label{lesson:JA-R142-second-pass-the-future-and-notices}



\begin{quote}
Three recalls, then the review.

\begin{itemize}
\raggedright
  \item \emph{Recall:} read \emph{no smoking} on a sign — \textbf{R2}, five lessons back
  \item \emph{Recall:} say \emph{is not}, plainly — \textbf{R3}, twenty lessons back
  \item \emph{Write it:} \textbf{\ja{け}} from memory — \textbf{R4}, eighty lessons back
\end{itemize}
\end{quote}

\subsection*{Your turn}
Say each one in Japanese.

\begin{itemize}
\raggedright
  \item \emph{Say it:} next week, next month, next year, a plan, probably
  \item \emph{Read it:} \textbf{\ja{ちゅうい}}, \textbf{\ja{きんえん}}, \textbf{\ja{うけつけ}}, \textbf{\ja{えいぎょうちゅう}}, \textbf{\ja{じゅんびちゅう}} — and say what each sign asks of you
\end{itemize}

\emph{Read it:} this notice on a clinic door, then say what you do

\begin{quote}
\textbf{\ja{じゅんびちゅう}}
\end{quote}

\begin{quote}
\textbf{\ja{うけつけで} \ja{なまえを} \ja{かいて} \ja{ください。}}
\end{quote}

\begin{quote}
\textbf{\ja{きんえん}}
\end{quote}

(It is not open yet; when it opens, write your name at reception; do not smoke.) Last, say one plan for next year with \textbf{\ja{つもり}}.

\begin{tcolorbox}[breakable,colback=teal!4,colframe=teal!35!black,title={Before you move on}]
Next week? (\textbf{\ja{らいしゅう}}, \emph{raishū}.) Next month? (\textbf{\ja{らいげつ}}, \emph{raigetsu}.) Next year? (\textbf{\ja{らいねん}}, \emph{rainen}.) A plan? (\textbf{\ja{つもり}}, \emph{tsumori}.) Probably? (\textbf{\ja{でしょう}}, \emph{deshō}.) Caution? (\textbf{\ja{ちゅうい}}, \emph{chūi}.) No smoking? (\textbf{\ja{きんえん}}, \emph{kin'en}.) Reception? (\textbf{\ja{うけつけ}}, \emph{uketsuke}.) Open? (\textbf{\ja{えいぎょうちゅう}}, \emph{eigyōchū}.) Not yet open? (\textbf{\ja{じゅんびちゅう}}, \emph{junbichū}.)
\end{tcolorbox}
